% =====================================================================
%  Informe del proyecto TEL252 — Firmas digitales basadas en hash
%
%  Compilar:  latexmk -pdf main.tex      (pdflatex + biber)
%  En Overleaf: Menu > Compiler: pdfLaTeX (biber se detecta solo).
%
%  Formato exigido por la rúbrica (Lección 0, ítem 3):
%    - máximo 7 páginas; portada con autores y título; índices;
%      introducción, marco teórico, metodología, resultados,
%      análisis y conclusión; citas y referencias APA;
%    - Arial/Helvetica 12, interlineado 1,15, texto justificado; PDF.
%
%  Notas pendientes: \pend{...} se ven en rojo mientras \borradortrue.
%  Para la entrega, cambiar a \borradorfalse y revisar que no quede nada.
% =====================================================================
\documentclass[12pt,letterpaper]{article}

% --- Modo borrador -----------------------------------------------------
\newif\ifborrador
\borradortrue   % <- cambiar a \borradorfalse para la entrega

% --- Idioma y tipografía ----------------------------------------------
\usepackage[T1]{fontenc}
\usepackage[utf8]{inputenc}
\usepackage[spanish,es-tabla,es-noshorthands]{babel}
\usepackage{helvet}                       % Helvetica (Nimbus Sans)
\renewcommand{\familydefault}{\sfdefault}
\usepackage[helvet]{sfmath}               % fórmulas en sans, coherente con el texto
\usepackage{setspace}
\setstretch{1.15}                         % interlineado 1,15
\usepackage[margin=2.5cm]{geometry}

% --- Paquetes ---------------------------------------------------------
\usepackage{amsmath}
\usepackage{graphicx}
\usepackage{booktabs}
\usepackage{subcaption}
\usepackage{xcolor}
\usepackage{enumitem}
\setlist{nosep}
\usepackage{csquotes}
\usepackage[style=apa,backend=biber]{biblatex}
\addbibresource{referencias.bib}
\usepackage[hidelinks]{hyperref}

% --- Macros del grupo -------------------------------------------------
% Nota pendiente para el grupo (desaparece con \borradorfalse).
\newcommand{\pend}[1]{\ifborrador\textcolor{red}{\textbf{[#1]}}\fi}

% Imagen de figuras/ o un recuadro si el archivo aún no existe.
% Uso: \grafico{archivo.pdf}
\newcommand{\grafico}[1]{%
  \IfFileExists{figuras/#1}%
    {\includegraphics[width=\linewidth]{figuras/#1}}%
    {\fbox{\parbox[c][3.5cm][c]{0.9\linewidth}{\centering\small
      Falta \texttt{figuras/\detokenize{#1}}\\(la genera \texttt{bench.py})}}}}

% Permite más figuras/tablas por página (evita páginas solo de flotantes).
\renewcommand{\topfraction}{0.9}
\renewcommand{\bottomfraction}{0.8}
\renewcommand{\textfraction}{0.07}
\renewcommand{\floatpagefraction}{0.8}

\newcommand{\B}{\,\mathrm{B}}   % bytes en fórmulas

% --- Datos del proyecto -----------------------------------------------
\newcommand{\titulo}{Firmas basadas en hash vs.\ ECDSA: el costo de resistir al cuántico}
\newcommand{\urlcodigo}{\url{https://github.com/josue-fourier/Tareas-TEL252-Crypto-2026-2/tree/main/Proyecto/codigo}}
\newcommand{\urlvideo}{\url{https://youtu.be/VIDEO}}             % <- reemplazar

\begin{document}

% =====================================================================
%  PORTADA
% =====================================================================
\begin{titlepage}
  \centering
  {\large Universidad Técnica Federico Santa María\\
   Departamento de Electrónica\par}
  \vspace{0.5cm}
  {TEL252 -- Criptografía y Seguridad de la Información\par}
  \vfill
  {\LARGE\bfseries \titulo\par}
  \vspace{1cm}
  {\large Proyecto grupal -- Equipo Hashquantum (grupo 12)\par}
  \vfill
  \begin{tabular}{ll}
    Integrante 1 & Josué N. Leiva Poblete -- 202430501-8 \\
    Integrante 2 & Sebastián Gabriel Moya Pérez -- \pend{Rol USM} \\
    Integrante 3 & Cristian Ignacio Pizarro Villaseñor -- \pend{Rol USM} \\
    Integrante 4 & Bastián Alexander Reyes Madrid -- \pend{Rol USM} \\
  \end{tabular}
  \vfill
  Profesora: Berioska Contreras Vargas\\[0.3cm]
  Santiago, noviembre de 2026
\end{titlepage}

% =====================================================================
%  ÍNDICES
% =====================================================================
\setcounter{tocdepth}{1}
\begingroup
  \setstretch{1.0}
  \tableofcontents
  \listoffigures
  \listoftables
\endgroup

% =====================================================================
\section{Introducción}
% Responsable: P4. Presupuesto: ~0,5 página.
% =====================================================================
Las firmas digitales garantizan autenticidad, integridad y no repudio, y hoy se basan principalmente en RSA, DSA y ECDSA \parencite{fips1865}. Su seguridad descansa en la dificultad de factorizar enteros y de calcular logaritmos discretos, dos problemas que el algoritmo de Shor resuelve en tiempo polinomial en un computador cuántico de escala suficiente \parencite{shor1994}. Por esta razón, el NIST propone dejar de permitir estas firmas después de 2035 \parencite{nistir8547}.

Las firmas basadas en hash ofrecen una alternativa cuya seguridad depende solo de la función hash, que el algoritmo de Grover debilita únicamente de forma cuadrática \parencite{grover1996}. El precio de esa seguridad es una firma más grande y, según los parámetros, más tiempo de cómputo. Este trabajo cuantifica ese precio: implementa desde cero tres esquemas de firma basados en hash y los compara con ECDSA P-256.

\begin{quote}
  \textbf{Hipótesis.} Una firma basada en hash (árbol de Merkle de altura $h=4$ con WOTS, $w=16$), implementada en Python puro, verifica en menos de 5 veces el tiempo de ECDSA P-256 implementado en C, a cambio de una firma entre 30 y 40 veces más grande. Además, pasar de $w=4$ a $w=256$ reduce la firma a menos de un tercio y multiplica por más de 10 el tiempo de firma y verificación.
\end{quote}

\textbf{Objetivo general.} Evaluar empíricamente el costo en tamaño y tiempo de una firma basada en hash implementada desde cero, frente a ECDSA P-256.

\textbf{Objetivos específicos.}
\begin{enumerate}
  \item Implementar las firmas de Lamport, WOTS y árbol de Merkle con SHA-256, y validarlas con pruebas de aceptación y rechazo.
  \item Medir el tamaño de claves y firmas y el tiempo de generar claves, firmar y verificar.
  \item Evaluar el efecto del parámetro $w$ sobre esas medidas.
  \item Comparar los resultados con ECDSA P-256 y demostrar con un ataque que reutilizar una clave de un solo uso permite falsificar firmas.
\end{enumerate}

% =====================================================================
\section{Marco teórico}
% Responsable: P3. Presupuesto: ~1,5 páginas.
% =====================================================================
\subsection{Funciones hash y firmas digitales}
Una función hash criptográfica $H$ transforma un mensaje de largo arbitrario en un resumen de largo fijo; SHA-256 produce 256 bits \parencite{fips1804}. Se le exige resistencia a preimagen (dado $y$, es inviable hallar $m$ con $H(m)=y$), a segunda preimagen y a colisiones \parencite{paar2010}. En un esquema de firma, quien posee la clave privada firma el resumen $H(m)$ y cualquiera con la clave pública verifica la firma. ECDSA sobre la curva P-256, el esquema de referencia de este trabajo, produce firmas de 64 bytes \parencite{fips1865}.

\subsection{La amenaza cuántica}
El algoritmo de Shor resuelve la factorización y el logaritmo discreto en tiempo polinomial, por lo que un computador cuántico suficientemente grande rompería RSA, DSA y ECDSA, y alargar las claves no lo evita \parencite{shor1994}. El algoritmo de Grover, en cambio, acelera la búsqueda exhaustiva solo de forma cuadrática: una preimagen de SHA-256 sigue costando del orden de $2^{128}$ operaciones \parencite{grover1996}. El borrador NIST IR 8547 clasifica ECDSA y RSA como vulnerables, propone dejar de permitirlas después de 2035 y nombra como reemplazos los esquemas basados en hash de FIPS 205 y SP 800-208 \parencite{nistir8547}.

\subsection{Firma de Lamport}
La clave privada son $2\times256$ valores aleatorios de 32 bytes, $sk_{i,b}$, y la clave pública es el hash de cada uno, $pk_{i,b}=H(sk_{i,b})$ \parencite{lamport1979}. Para firmar $m$ se revela $sk_{i,b_i}$, donde $b_i$ es el bit $i$ de $H(m)$; para verificar se comprueba que $H(sk_{i,b_i})=pk_{i,b_i}$. Falsificar exige invertir el hash. La firma mide $8\,192$ bytes y las claves $16\,384$ bytes cada una. La clave es de un solo uso: cada firma revela la mitad de los secretos.

\subsection{Firma de Winternitz (WOTS)}
WOTS reduce la firma firmando varios bits a la vez \parencite{merkle1989}. El resumen se divide en $\ell_1$ dígitos en base $w$, y a cada dígito se le asocia una cadena que parte de un secreto $x$ y aplica $H$ hasta $w-1$ veces; el último valor, $H^{w-1}(x)$, es la clave pública de la cadena. Para firmar el dígito $d$ se revela $H^{d}(x)$, y para verificar se aplica $H$ otras $w-1-d$ veces. Una suma de control $C=\sum_i (w-1-d_i)$, firmada en $\ell_2$ dígitos extra, impide que un atacante avance en las cadenas: subir un dígito obliga a bajar otro, lo que exigiría invertir $H$. Con resúmenes de 256 bits:
\begin{equation}
  \ell_1=\left\lceil \frac{256}{\log_2 w}\right\rceil,\qquad
  \ell_2=\left\lfloor \frac{\log_2\!\big(\ell_1 (w-1)\big)}{\log_2 w}\right\rfloor+1,\qquad
  |\sigma_{\mathrm{WOTS}}|=32\,(\ell_1+\ell_2)\B .
\end{equation}
Al aumentar $w$ baja el número de cadenas y crece su largo: la firma se achica y el cómputo aumenta.

\subsection{Árbol de Merkle}
Para firmar más de un mensaje con una sola clave pública se generan $2^h$ claves WOTS y se combinan en un árbol binario de hash: cada hoja es el hash de una clave pública WOTS y cada nodo es el hash de la concatenación de sus hijos \parencite{merkle1989}. La raíz, de 32 bytes, es la clave pública. La firma del mensaje $i$ contiene la firma WOTS con la clave $i$, la clave pública WOTS comprimida (32 bytes), los $h$ nodos hermanos del camino hasta la raíz y el índice:
\begin{equation}
  |\sigma_{\mathrm{Merkle}}| = |\sigma_{\mathrm{WOTS}}| + 32 + 32h + 4\B .
\end{equation}
Incluir la clave pública WOTS es redundante, porque el verificador la reconstruye desde la firma; cuesta 32 bytes y simplifica la verificación.
El esquema tiene estado: el firmante debe recordar qué índices usó y no repetirlos nunca. XMSS y LMS, estandarizados por el NIST en SP 800-208, siguen este diseño \parencite{nist800208}. SLH-DSA (FIPS 205) elimina el estado a cambio de firmas más grandes \parencite{fips205}. \pend{P3: confirmar en FIPS 205 los tamaños de firma del nivel de 128 bits}

% =====================================================================
\section{Metodología}
% Responsable: P4. Presupuesto: ~1 página.
% =====================================================================
La investigación es experimental y se desarrolla en cuatro fases.

\begin{enumerate}
  \item \textbf{Implementación.} Lamport, WOTS ($w\in\{4,16,256\}$) y árbol de Merkle ($h=4$) en Python 3.13.16, usando solo \texttt{hashlib} (SHA-256) a través de una función $H$ común. Los tres esquemas comparten la interfaz \texttt{keygen}, \texttt{sign} y \texttt{verify}. ECDSA P-256 se toma de la librería \texttt{cryptography} y se usa solo como referencia. El código está en \urlcodigo.
  \item \textbf{Validación.} Cada esquema pasa tres pruebas: una firma válida verifica, un mensaje alterado no verifica y una firma hecha con otra clave no verifica.
  \item \textbf{Medición.} Cada operación se repite $N=10\,000$ veces con \texttt{time.perf\_counter}, descartando 10 ejecuciones previas de calentamiento, y se reporta media y desviación estándar. Los tamaños se obtienen del largo en bytes de claves y firmas; para ECDSA se usa la firma en formato $r\Vert s$ (64 bytes). Todo se mide en el mismo computador: un AMD Ryzen 5 5600GT con 14 GiB de RAM y Ubuntu 26.04 LTS, dentro de un contenedor Docker de JupyterLab.
  \item \textbf{Análisis.} Se contrasta cada parte de la hipótesis con su umbral; se compara el tiempo medido con el costo teórico en hashes; se discute la validez de comparar Python con C; y se declaran las limitaciones.
\end{enumerate}

La variable independiente es $w$. Las dependientes son el tamaño de la clave pública y de la firma, y los tiempos de generación de claves, firma y verificación. Se controlan la función hash, la altura $h=4$, el número de repeticiones, la versión de Python y el computador.

\textbf{Prueba de concepto.} La demostración, grabada en el video (\urlvideo), tiene tres partes. Primero, cada esquema firma, verifica y rechaza un mensaje alterado. Segundo, se firman $k=6$ mensajes con una misma clave Lamport; el atacante busca un mensaje de la forma \enquote{Transferir \$1000 a Mallory \#$n$} cuyo resumen coincida con los secretos ya revelados. Con $k$ firmas, la cantidad esperada de bits que el atacante no puede elegir es $256/2^{k-1}$, es decir, 8 bits y unos $2^{8}$ intentos para $k=6$, frente a $2^{128}$ con solo dos firmas. Tercero, el árbol de Merkle firma $2^h=16$ mensajes y el firmante rechaza el decimoséptimo. La PoC es exitosa si las tres pruebas pasan en los tres esquemas, el ataque produce una firma aceptada y el firmante de Merkle rechaza el índice agotado.

% =====================================================================
\section{Resultados}
% Responsable: P2. Presupuesto: ~1 página.
% =====================================================================
Con $w=16$, la firma del árbol de Merkle mide 2\,308 bytes, 36,1 veces la de ECDSA P-256, y se verifica en 0,163 ms, 2,5 veces el tiempo de ECDSA (Tablas~\ref{tab:tamanos} y~\ref{tab:tiempos}, Figura~\ref{fig:resultados}).

\begin{table}[htbp]
  \centering
  \caption{Tamaño de claves y firmas: valor teórico y medido.}
  \label{tab:tamanos}
  \begin{tabular}{lrrr}
    \toprule
    Esquema & Clave pública (B) & Firma teórica (B) & Firma medida (B) \\
    \midrule
    Lamport                  & 16\,384 & 8\,192 & \pend{--} \\
    WOTS, $w=16$             & 2\,144  & 2\,144 & 2\,144 \\
    Merkle $h=4$, $w=4$      & 32      & 4\,420 & 4\,420 \\
    Merkle $h=4$, $w=16$     & 32      & 2\,308 & 2\,308 \\
    Merkle $h=4$, $w=256$    & 32      & 1\,252 & 1\,252 \\
    ECDSA P-256              & 64      & 64     & 64 \\
    \bottomrule
  \end{tabular}
\end{table}

\begin{table}[htbp]
  \centering
  \caption{Tiempo por operación en milisegundos (media $\pm$ desviación estándar, $N=10\,000$).}
  \label{tab:tiempos}
  \small
  \begin{tabular}{lrrrr}
    \toprule
    Esquema & Gen. claves & Firmar & Verificar & Verif./ECDSA \\
    \midrule
    Merkle $h=4$, $w=4$   & 3,415 $\pm$ 0,065 & 0,087 $\pm$ 0,005 & 0,098 $\pm$ 0,004 & 1,5 \\
    Merkle $h=4$, $w=16$  & 5,251 $\pm$ 0,093 & 0,158 $\pm$ 0,006 & 0,163 $\pm$ 0,008 & 2,5 \\
    Merkle $h=4$, $w=256$ & 38,168 $\pm$ 0,631 & 1,124 $\pm$ 0,029 & 1,270 $\pm$ 0,033 & 19,7 \\
    ECDSA P-256           & 0,014 $\pm$ 0,001 & 0,032 $\pm$ 0,002 & 0,064 $\pm$ 0,005 & 1 \\
    \bottomrule
  \end{tabular}
\end{table}

\begin{figure}[htbp]
  \centering
  \begin{subfigure}[t]{0.48\linewidth}
    \grafico{tamano_tiempo_w.pdf}
    \caption{Efecto de $w$: tamaño de firma y tiempos de firma y verificación (hipótesis 2).}
    \label{fig:w}
  \end{subfigure}\hfill
  \begin{subfigure}[t]{0.48\linewidth}
    \grafico{comparacion_ecdsa.pdf}
    \caption{Razón frente a ECDSA P-256 en tamaño y en verificación (hipótesis 1).}
    \label{fig:ecdsa}
  \end{subfigure}
  \caption{Resultados de la medición (Merkle, $h=4$; $N=10\,000$). Líneas y bandas: umbrales de la hipótesis.}
  \label{fig:resultados}
\end{figure}

\textbf{Prueba de concepto.} En el árbol de Merkle ($h=4$, $w=16$) las 16 firmas verificaron; la firma de un mensaje alterado y una firma hecha con otro árbol fueron rechazadas; y el firmante rechazó el decimoséptimo mensaje con el error de claves agotadas. \pend{P1: resultado de las tres pruebas para Lamport y WOTS; en el ataque, bits fijos y cantidad de intentos observados frente a los esperados ($\approx 2^8$).}

% =====================================================================
\section{Análisis}
% Responsables: P2 y P4. Presupuesto: ~0,75 página.
% =====================================================================
\textbf{Contraste con la hipótesis.} Las cuatro partes de la hipótesis se cumplen (Figura~\ref{fig:resultados}). Con $w=16$, verificar toma 2,5 veces el tiempo de ECDSA (umbral: 5) y la firma es 36,1 veces más grande (rango esperado: 30 a 40). Al pasar de $w=4$ a $w=256$, la firma se reduce a 0,28 veces su tamaño (umbral: menos de un tercio), mientras que firmar se vuelve 12,9 veces más lento y verificar, 12,9 veces (umbral: más de 10).

\textbf{Predicción frente a medición.} En promedio, firmar con WOTS cuesta $\ell\,(w-1)/2$ evaluaciones de $H$, y verificar otras tantas. En el árbol, verificar suma $1+h$ hashes: la hoja y el camino hasta la raíz. Así, se esperan 204,5, 507,5 y 4\,340 evaluaciones para $w=4$, 16 y 256 ($\ell=133$, 67 y 34), un aumento de 21,2 veces. El tiempo medido crece menos, 12,9 veces (línea punteada de la Figura~\ref{fig:w}), porque el costo por hash baja de 0,48 a 0,29~$\mu$s. En Python, cada cadena agrega un costo fijo del intérprete (llamadas, bucles y cortes de bytes) que pesa más cuando hay muchas cadenas cortas, como con $w=4$. Por lo tanto, el número de hashes predice bien el orden de magnitud, pero sobrestima el efecto de $w$.

\textbf{Validez de la comparación.} ECDSA se ejecuta en código C optimizado, mientras que los esquemas propios corren en Python puro; la comparación favorece a ECDSA, de modo que una diferencia pequeña es un resultado conservador. Con $N=10\,000$, la desviación estándar de los esquemas de Merkle fue de entre 2 y 5~\% de la media, y la de ECDSA de entre 6 y 10~\%. Las diferencias entre esquemas son mucho mayores que esa dispersión, por lo que las conclusiones no dependen del ruido de medición.

\textbf{Limitaciones y contexto.} La implementación tiene estado, no es de tiempo constante y no incluye SLH-DSA ni W-OTS+. \pend{Relacionar con la migración que propone el NIST: qué significa el tamaño medido frente a ECDSA y frente a SLH-DSA.}

% =====================================================================
\section{Conclusión}
% Responsables: P2 y P4. Presupuesto: ~0,25 página.
% =====================================================================
\pend{Responder el problema en 3 a 5 frases: cuánto cuesta en tamaño y tiempo una firma basada en hash frente a ECDSA, qué partes de la hipótesis se cumplieron, qué demuestra el ataque a Lamport y qué quedaría como trabajo futuro (W-OTS+, SLH-DSA, implementación en C).}

% =====================================================================
{\sloppy\printbibliography[title={Referencias}]}
% =====================================================================

% =====================================================================
\appendix
\section{Retroalimentación entre pares}
% Responsable: P4. Ítem 10 de la rúbrica: incluir la dada y la recibida.
% =====================================================================
\textbf{Retroalimentación dada} por el equipo \pend{nuestro equipo} al equipo \pend{equipo revisado}:

\noindent\begin{tabular}{p{0.25\linewidth}p{0.68\linewidth}}
  \toprule
  Aspecto & Comentario \\
  \midrule
  \pend{p.\,ej., hipótesis} & \pend{comentario entregado} \\
  \bottomrule
\end{tabular}

\vspace{0.5cm}
\textbf{Retroalimentación recibida} del equipo \pend{equipo revisor}:

\noindent\begin{tabular}{p{0.25\linewidth}p{0.4\linewidth}p{0.25\linewidth}}
  \toprule
  Aspecto & Comentario & Cómo se incorporó \\
  \midrule
  \pend{aspecto} & \pend{comentario recibido} & \pend{cambio hecho} \\
  \bottomrule
\end{tabular}

\end{document}
